The model must be well-defined to be used in NMPC, in such a way 
to obtain reliable predictions of system behavior in response to various control actions.
The usage of a data-driven approach, such as neural networks, in the vehicle modeling phase has the
advantage of not modeling all the nonlinear physical effects, which can be also very complex.
In fact, neural networks have the capability to accurately model them. 
Neural networks are universal approximators, and they can be used both for designing 
the entire dynamic model of a system and approximating the disturbances and noises 
a system can be affected.
Our purpose is to approximate the entire kinematic model of an unicycle with a neural model,
due to its low complexity. 
In this section we will start describing briefly the kinematic model of a unicycle, and then 
we will describe the neural architecture used to approximate it.

\subsection{Unicycle kinematic model}
The kinematic model of the unicyle is a simplified representation of the movement of the vehicle
in temrs of position and orientation. It assumes there will be no lateral deformations. 
The state space is represented by its position and orientational
\begin{equation}
    [x, y, \theta]
\end{equation} 
while the input actions are represented by longitudinal and angular velocities the vehicle can 
assume
\begin{equation}
    [v, \omega]
\end{equation}
So, the motion of the unicycle can be described through the following equations:
\begin{eqnarray}
    \dot{x}(t) &=& v(t) \cos(\theta(t)) \\
    \dot{y}(t) &=& v(t) \sin(\theta(t)) \\
    \dot{\theta}(t) &=& \omega(t) \\
\end{eqnarray}

\subsection{Unicycle Neural model}

To model the kinematic model of the unicycle, a simple architecture is
needed, but not too simple, in order to captivate the dynamics of the
system.

\subsubsection{Architecture}\label{architecture} The neural network used to
model the kinematic model of the unicycle is a simple MLP architecture,
with two hidden layers of 128 neurons. The network has two inputs which are
concatenated together, the current state of the unicycle (dimension 3) and
the current applied actions on the vehicle (dimension 2), so it must be of
dimension 5 and is represented by 
\begin{equation}
    X = [x, y, \theta, v, \omega]
\end{equation}

As for the output, we tried two different strategies: the first one is to
have the output of the network to be the next state of the unicycle, so it
is a vector of dimension 3, which represents the position and orientation
of the vehicle at the next time-step. The second strategy is to have the
output of the network to be the dynamics of the system, so it is a vector
of dimension 3, which represents the velocity of the vehicle in the x-y
direction and the angular velocity.

In the first case, the output at time $k$ is
\begin{equation}
    Y_k = [x_{k+1}, y_{k+1}, \theta_{k+1}]
\end{equation}

In the second case, the output at time $k$ is
\begin{equation}
    Y_k = [\dot{x}, \dot{y}_k, \dot{\theta}_k]
\end{equation}

\subsubsection{Dataset}
The dataset was synthetically generated by using the nominal model of the
robot and generating $N$ trajectories, each composed by $K$ steps, where
the inputs are maintained constant for a number of $M$ steps. This choice was
taken because it was noticed that the interval between consecutive
time-steps is too small, therefore the changes about the state of the
unicycle are small too, and the network was not able to approximate
correctly the motion of the vehicle. On the contrary, having the same input
for multiple steps allows to understand in a better way the dynamics of the
system subject to specific inputs. \\
Each sample is not just a pair state-action at step $t$, but it is a window
of dimension $k$ which contains all pairs state-action from time-step $t$
till $t+k$ (less the last action $a_{t+k}$ which is useless), so it is in
the form 
\begin{equation}
    (s_t, a_t, s_{t+1}, a_{t+1}, ..., s_{t+k-1}, a_{t+k-1}, s_{t+k})
\end{equation}
This allows to have a larger window of evolution of the system with respect
to just the motion executed in one single step. What really matters, of
course, are the labels of our dataset. They are either the derivatives, or
the successive states. 

An examplary configuration of the dataset is the following:
\begin{itemize}
    \item $N = 10.000$
    \item $K = 20$
    \item $M = 5$
\end{itemize}

Each sub-trajectory ends with the first state of the successive
sub-trajectory. This way we are able to learn the dynamics of the system
under different time conditions, in order to observe better the changes and
try the network to learn better.

We divided the dataset in training, validation and test set, such that the
validation set can guide the training process, while the test set can
validate the results, in thi case also graphically.


\subsubsection{Training and loss}

As said before, we tried two different strategies to train the network. In
the first case, the network has to learn the derivative of the system,
given a state and an input. Thus, we simply compare the output of the
network with the nominal model, imposing a mean squared error loss:
\begin{equation}
    loss = \frac{1}{k}\sum_{i=t}^{t+k} (\bar{s}_i - s_i)^2
\end{equation}

In the second case, the network has to learn the next state of the system,
thus we still use a mean squared error loss, but this time on the states.


\subsubsection{Extra: PINN}

We also tried to make our network a Physics-Informed Neural network, which
incorporates some prior physical knowledge through the addition of a
specific term in the loss function. In detail, the neural architecture is
the same explained in \autoref{architecture}; what changes is that, if we
have already some information about the dynamics of the system in the form
of differential equations, we can add a physics loss term telling how much
the approximated dynamics reflects the known one. In this way, the loss
function becomes:
\begin{equation}
    loss = \frac{1}{k}\sum_{i=t}^{t+k} (\bar{s}_i - s_i)^2 + \lambda \frac{1}{M}\sum_{j}^{M}(g(x_j) - f(x_j))^2
\end{equation}

where $g$ is the known differential equation, while $f$ is the neural approximator. The $M$ points
on which the physics loss is computed are called \textbf{collocation points}. These must not be
samples in the training datasets, but they can be any points we want, of which we can compute the 
dynamics. This can be a useful approach if we don't have a large amount of training data and if we
already know something about the problem which is not related directly to data. 